% =============================================================================
% Documentation — Intégration de HashiCorp Vault dans le pipeline CI/CD Jenkins
% Compilation : pdflatex vault-documentation.tex   (2 passes pour la table des matières)
%               ou import direct sur Overleaf
% =============================================================================
\documentclass[11pt,a4paper]{article}

\usepackage[utf8]{inputenc}
\usepackage[T1]{fontenc}
\usepackage{lmodern}
\usepackage[french]{babel}
\usepackage[margin=2.3cm]{geometry}
\usepackage{amssymb}
\usepackage{xcolor}
\usepackage{listings}
\usepackage{booktabs}
\usepackage{tabularx}
\usepackage{array}
\usepackage{enumitem}
\usepackage{fancyhdr}
\usepackage[most]{tcolorbox}
\usepackage{tikz}
\usetikzlibrary{positioning,arrows.meta,shapes.geometric,babel}
\usepackage[hidelinks]{hyperref}

% --- Couleurs -----------------------------------------------------------------
\definecolor{vaultyellow}{HTML}{FFD814}
\definecolor{vaultdark}{HTML}{1F2937}
\definecolor{codebg}{HTML}{F6F8FA}
\definecolor{codeframe}{HTML}{D0D7DE}
\definecolor{kw}{HTML}{0550AE}
\definecolor{str}{HTML}{0A3069}
\definecolor{cmt}{HTML}{6E7781}
\definecolor{okgreen}{HTML}{1A7F37}
\definecolor{warnred}{HTML}{CF222E}

% --- Listings -----------------------------------------------------------------
\lstset{
  basicstyle=\ttfamily\small,
  backgroundcolor=\color{codebg},
  frame=single, rulecolor=\color{codeframe},
  keywordstyle=\color{kw}\bfseries,
  stringstyle=\color{str},
  commentstyle=\color{cmt}\itshape,
  breaklines=true, breakatwhitespace=true,
  columns=fullflexible, keepspaces=true,
  showstringspaces=false,
  xleftmargin=4pt, xrightmargin=4pt,
  aboveskip=8pt, belowskip=8pt,
  literate=
    {é}{{\'e}}1 {è}{{\`e}}1 {ê}{{\^e}}1 {à}{{\`a}}1 {â}{{\^a}}1
    {ç}{{\c{c}}}1 {ù}{{\`u}}1 {ô}{{\^o}}1 {î}{{\^i}}1 {É}{{\'E}}1
    {→}{{$\rightarrow$}}1 {⚠}{{!}}1
}
\lstdefinelanguage{yaml}{
  keywords={true,false,null},
  sensitive=true,
  comment=[l]{\#},
  morestring=[b]",
  morestring=[b]'
}
\lstdefinelanguage{hcl}{
  keywords={path,capabilities},
  comment=[l]{\#},
  morestring=[b]"
}
\lstdefinelanguage{groovy}{
  keywords={def,new,return,import,if,else,try,catch,String,withVault,sh},
  sensitive=true,
  comment=[l]{//},
  morecomment=[s]{/*}{*/},
  morestring=[b]",
  morestring=[b]'
}

% --- Encadrés -----------------------------------------------------------------
\newtcolorbox{note}[1][Note]{colback=blue!4,colframe=kw,fonttitle=\bfseries,title=#1}
\newtcolorbox{attention}[1][Attention]{colback=red!4,colframe=warnred,fonttitle=\bfseries,title=#1}
\newtcolorbox{retenir}[1][À retenir]{colback=green!4,colframe=okgreen,fonttitle=\bfseries,title=#1}

% --- En-têtes -----------------------------------------------------------------
\pagestyle{fancy}
\fancyhf{}
\fancyhead[L]{\small HashiCorp Vault — Pipeline CI/CD}
\fancyhead[R]{\small Demo Sécurité App Moderne}
\fancyfoot[C]{\thepage}

\newcommand{\vpath}[1]{\texttt{#1}}

\title{\textbf{Intégration de HashiCorp Vault\\dans le pipeline CI/CD Jenkins}\\[0.4em]
       \large Gestion centralisée des secrets — Guide technique}
\author{Demo Sécurité App Moderne}
\date{\today}

% =============================================================================
\begin{document}

\maketitle
\thispagestyle{empty}

\begin{abstract}
Ce document explique pourquoi et comment HashiCorp Vault a été ajouté à la chaîne
CI/CD (Jenkins, SonarQube, OWASP Dependency-Check, Trivy, Docker Hub). Les secrets du
pipeline (token Docker Hub, clé NVD, secrets clients Keycloak) ne sont plus saisis
manuellement dans Jenkins : ils sont stockés dans Vault et lus par le pipeline au
moment du build, via une identité machine (\emph{AppRole}) limitée par une
\emph{policy} en lecture seule. Le document couvre les concepts, l'architecture, chaque
fichier modifié, le flux d'exécution, l'exploitation et les différences entre le mode
démonstration et une installation de production.
\end{abstract}

\tableofcontents
\newpage

% =============================================================================
\section{Contexte et objectifs}
% =============================================================================

\subsection{Le problème des secrets dans un pipeline}

Un pipeline CI/CD manipule de nombreux secrets : identifiants de registre Docker,
clés d'API, secrets OAuth2/OIDC, mots de passe de bases de données\ldots{} Avant
l'intégration de Vault, ces secrets étaient enregistrés à la main dans
\emph{Jenkins Credentials}, ce qui pose plusieurs problèmes :

\begin{itemize}[leftmargin=*]
  \item \textbf{Dispersion} : chaque outil (Jenkins, applications, scripts) possède sa
        propre copie des secrets ;
  \item \textbf{Traçabilité faible} : difficile de savoir qui a lu quel secret et quand ;
  \item \textbf{Rotation pénible} : changer un mot de passe impose de le modifier partout ;
  \item \textbf{Reproductibilité} : à chaque réinstallation de Jenkins, il faut ressaisir
        les credentials manuellement ;
  \item \textbf{Risque de fuite} : tentation de mettre les secrets dans le code, les
        fichiers \texttt{docker-compose} ou les variables d'environnement en clair.
\end{itemize}

\subsection{Objectifs}

\begin{enumerate}[leftmargin=*]
  \item \textbf{Centraliser} tous les secrets du pipeline dans un coffre-fort unique ;
  \item \textbf{Authentifier} Jenkins comme une machine, sans mot de passe dans le code ;
  \item Appliquer le \textbf{principe du moindre privilège} : Jenkins ne peut que
        \emph{lire} les secrets dont il a besoin ;
  \item \textbf{Automatiser} entièrement la configuration (aucune action manuelle dans
        Jenkins) ;
  \item Ne \textbf{jamais} versionner de secret dans Git.
\end{enumerate}

% =============================================================================
\section{Concepts clés de Vault}
% =============================================================================

\begin{description}[leftmargin=0pt, style=nextline]
  \item[Secret engine KV v2]
    Moteur de stockage clé/valeur versionné. En mode dev, il est monté sur
    \vpath{secret/}. Un secret est un ensemble de paires clé/valeur à un chemin donné,
    par exemple \vpath{secret/ci/dockerhub} contient \texttt{username} et
    \texttt{password}. En KV v2, l'API insère \vpath{data/} dans le chemin :
    \vpath{secret/data/ci/dockerhub}.

  \item[Méthode d'authentification (auth method)]
    Manière dont un client prouve son identité à Vault : token, userpass, LDAP, OIDC,
    Kubernetes, \textbf{AppRole}\ldots

  \item[AppRole]
    Méthode d'authentification conçue pour les \textbf{machines et applications}.
    Elle repose sur deux éléments :
    \begin{itemize}
      \item \texttt{role\_id} : identifiant du rôle (équivalent d'un nom d'utilisateur,
            non secret) ;
      \item \texttt{secret\_id} : secret associé (équivalent d'un mot de passe).
    \end{itemize}

  \item[Policy]
    Règles écrites en HCL qui définissent, chemin par chemin, les opérations autorisées
    (\texttt{read}, \texttt{create}, \texttt{update}, \texttt{delete}, \texttt{list}).
    Tout ce qui n'est pas explicitement autorisé est \textbf{refusé}.

  \item[Token]
    Après authentification, Vault délivre un token \textbf{temporaire} (TTL) auquel sont
    attachées les policies. C'est ce token qui est utilisé pour lire les secrets.

  \item[Seal / Unseal]
    Au démarrage, Vault est \emph{scellé} : ses données sont chiffrées et il ne peut pas
    les lire. Il doit être \emph{descellé} (clés d'unseal ou auto-unseal via un KMS).
    Le mode dev démarre directement descellé.
\end{description}

% =============================================================================
\section{Architecture}
% =============================================================================

\subsection{Vue d'ensemble}

\begin{center}
\begin{tikzpicture}[
  node distance=1.2cm and 1.5cm,
  box/.style={draw, rounded corners, thick, minimum width=2.9cm, minimum height=1.1cm,
              align=center, font=\small},
  file/.style={box, fill=codebg, draw=codeframe},
  vault/.style={box, fill=vaultyellow!70, draw=vaultdark},
  svc/.style={box, fill=blue!8, draw=kw},
  arr/.style={-{Stealth[length=3mm]}, thick},
  lbl/.style={font=\scriptsize, align=center, fill=white, inner sep=1.5pt}
]
  \node[file]  (env)   {\texttt{secrets.env}\\\scriptsize(local, ignoré par Git)};
  \node[svc,   right=of env]   (init)  {\textbf{vault-init}\\\scriptsize(one-shot)};
  \node[vault, right=of init]  (vault) {\textbf{Vault}\\\scriptsize KV v2 + AppRole};
  \node[file,  below=1.6cm of init]  (vol)  {volume\\\texttt{vault\_approle}};
  \node[svc,   below=1.6cm of vault] (jenk) {\textbf{Jenkins}\\\scriptsize plugin Vault};
  \node[box,   right=of jenk, fill=green!6, draw=okgreen] (pipe) {Pipeline\\\texttt{withVault}};

  \draw[arr] (env) -- node[lbl, above]{valeurs} (init);
  \draw[arr] (init) -- node[lbl, above]{configure\\(token root)} (vault);
  \draw[arr] (init) -- node[lbl, left]{écrit\\role\_id / secret\_id} (vol);
  \draw[arr] (vol) -- node[lbl, below]{lu au démarrage} (jenk);
  \draw[arr] (jenk) -- node[lbl, right]{AppRole\\login} (vault);
  \draw[arr] (jenk) -- node[lbl, below]{secrets} (pipe);
\end{tikzpicture}
\end{center}

\subsection{Ordre de démarrage des conteneurs}

\begin{center}
\begin{tabularx}{\linewidth}{c l X}
\toprule
\textbf{Ordre} & \textbf{Conteneur} & \textbf{Rôle} \\
\midrule
1 & \texttt{vault}      & Serveur Vault (mode dev), port 8200 \\
2 & \texttt{vault-init} & Configure Vault (AppRole, policy, secrets) puis s'arrête \\
3 & \texttt{jenkins}    & Démarre uniquement si \texttt{vault-init} s'est terminé avec succès \\
\bottomrule
\end{tabularx}
\end{center}

Cet ordre est garanti par les conditions \texttt{depends\_on} de Docker Compose :
\texttt{vault-init} attend que \texttt{vault} soit \emph{healthy}, et \texttt{jenkins}
attend que \texttt{vault-init} soit \emph{completed\_successfully}.

\subsection{Fichiers concernés}

\begin{center}
\begin{tabularx}{\linewidth}{>{\ttfamily\small}l X}
\toprule
\textrm{\textbf{Fichier}} & \textbf{Rôle} \\
\midrule
ci/docker-compose-ci-tools.yml          & Ajout des services \texttt{vault} et \texttt{vault-init}, du volume partagé \\
ci/vault/init.sh                        & Script d'initialisation de Vault \\
ci/vault/jenkins-policy.hcl             & Policy lecture seule pour Jenkins \\
ci/vault/secrets.env.example            & Modèle des valeurs de secrets \\
ci/jenkins/plugins.txt                  & Ajout du plugin \texttt{hashicorp-vault-plugin} \\
ci/jenkins/init.groovy.d/03-vault.groovy & Configuration automatique du plugin dans Jenkins \\
Jenkinsfile                             & \texttt{withCredentials} remplacé par \texttt{withVault} \\
.gitignore                              & Exclusion de \texttt{ci/vault/secrets.env} \\
.gitattributes                          & Fins de ligne LF forcées pour les scripts \texttt{*.sh} \\
\bottomrule
\end{tabularx}
\end{center}

% =============================================================================
\section{Mise en œuvre détaillée}
% =============================================================================

\subsection{Étape 1 — Le serveur Vault}

Service ajouté dans \vpath{ci/docker-compose-ci-tools.yml} :

\begin{lstlisting}[language=yaml]
vault:
  image: hashicorp/vault:1.17
  container_name: vault
  command: server -dev
  ports:
    - "8200:8200"
  environment:
    VAULT_DEV_ROOT_TOKEN_ID: ${VAULT_ROOT_TOKEN:-root}
    VAULT_DEV_LISTEN_ADDRESS: 0.0.0.0:8200
  cap_add:
    - IPC_LOCK
  networks: [ci-net]
  healthcheck:
    test: ["CMD-SHELL", "wget -qO- http://127.0.0.1:8200/v1/sys/health || exit 1"]
\end{lstlisting}

\begin{itemize}[leftmargin=*]
  \item \texttt{server -dev} : mode développement, stockage en mémoire, déjà descellé,
        token administrateur fixe ;
  \item \texttt{VAULT\_DEV\_ROOT\_TOKEN\_ID} : valeur du token root (\texttt{root} par
        défaut, surchargeable par la variable \texttt{VAULT\_ROOT\_TOKEN}) ;
  \item \texttt{cap\_add: IPC\_LOCK} : autorise Vault à verrouiller sa mémoire pour
        empêcher que des secrets soient écrits dans le \emph{swap} ;
  \item \texttt{networks: [ci-net]} : Jenkins joint Vault par son nom DNS
        \texttt{http://vault:8200}.
\end{itemize}

\subsection{Étape 2 — L'initialisation automatique (\texttt{vault-init})}

\begin{lstlisting}[language=yaml]
vault-init:
  image: hashicorp/vault:1.17
  user: root
  entrypoint: ["/bin/sh", "/vault/config/init.sh"]
  environment:
    VAULT_ADDR: http://vault:8200
    VAULT_TOKEN: ${VAULT_ROOT_TOKEN:-root}
  env_file:
    - path: ./vault/secrets.env
      required: false
  volumes:
    - ./vault/init.sh:/vault/config/init.sh:ro
    - ./vault/jenkins-policy.hcl:/vault/config/jenkins-policy.hcl:ro
    - vault_approle:/vault/approle
  depends_on:
    vault:
      condition: service_healthy
  restart: "no"
\end{lstlisting}

Ce conteneur réutilise l'image Vault (qui contient la CLI \texttt{vault}), exécute
\vpath{init.sh} une seule fois, puis s'arrête. Le script est \textbf{idempotent} : il peut
être relancé à chaque \texttt{docker compose up} sans effet indésirable.

\subsubsection*{2.a — Création de l'AppRole Jenkins}

\begin{lstlisting}[language=bash]
vault auth list | grep -q '^approle/' || vault auth enable approle
vault policy write jenkins-ci /vault/config/jenkins-policy.hcl
vault write auth/approle/role/jenkins-ci \
    token_policies=jenkins-ci token_ttl=1h token_max_ttl=4h secret_id_ttl=0
vault write auth/approle/role/jenkins-ci/role-id role_id=jenkins-ci
\end{lstlisting}

\begin{itemize}[leftmargin=*]
  \item active la méthode AppRole si elle ne l'est pas déjà ;
  \item enregistre la policy \texttt{jenkins-ci} ;
  \item crée le rôle : les tokens délivrés porteront la policy \texttt{jenkins-ci},
        vivront 1\,h (4\,h maximum) ;
  \item fixe un \texttt{role\_id} lisible et stable : \texttt{jenkins-ci}.
\end{itemize}

\subsubsection*{2.b — Génération d'un \texttt{secret\_id} stable}

\begin{lstlisting}[language=bash]
[ -s "$APPROLE_DIR/secret_id" ] || \
    cat /proc/sys/kernel/random/uuid > "$APPROLE_DIR/secret_id"
echo "$ROLE" > "$APPROLE_DIR/role_id"

SECRET_ID=$(cat "$APPROLE_DIR/secret_id")
vault write "auth/approle/role/$ROLE/custom-secret-id" secret_id="$SECRET_ID"
# Vérifie que Jenkins pourra réellement s'authentifier
vault write auth/approle/login role_id="$ROLE" secret_id="$SECRET_ID"
\end{lstlisting}

Le \texttt{secret\_id} est généré aléatoirement \textbf{une seule fois}, puis stocké dans
le volume Docker \texttt{vault\_approle}. À chaque redémarrage, il est ré-enregistré dans
Vault : Jenkins garde ainsi des identifiants valides même si Vault (en mémoire) a été
réinitialisé. Le login de test final fait échouer \texttt{vault-init} (et donc bloque le
démarrage de Jenkins) si l'authentification ne fonctionne pas.

\begin{retenir}
Le \texttt{secret\_id} n'apparaît \textbf{jamais} dans le code source ni dans Git : il
transite uniquement par un volume Docker partagé entre \texttt{vault-init} (écriture)
et \texttt{jenkins} (lecture seule).
\end{retenir}

\subsubsection*{2.c — Création des secrets}

\begin{lstlisting}[language=bash]
seed() {
    path=$1; shift
    if vault kv get "$path" >/dev/null 2>&1; then
        echo "[vault-init] $path existe déjà — conservé"
    else
        vault kv put "$path" "$@" >/dev/null
        echo "[vault-init] $path créé"
    fi
}

seed secret/ci/dockerhub \
    username="${DOCKERHUB_USERNAME:-moustaphisene}" \
    password="${DOCKERHUB_TOKEN:-CHANGE_ME}"
seed secret/ci/nvd api_key="${NVD_API_KEY:-CHANGE_ME}"
\end{lstlisting}

La fonction \texttt{seed} ne crée un secret que s'il n'existe pas : une valeur modifiée
dans l'interface de Vault n'est jamais écrasée par une relance du script. Les valeurs
proviennent du fichier local \vpath{ci/vault/secrets.env}.

\subsection{Étape 3 — La policy (moindre privilège)}

Fichier \vpath{ci/vault/jenkins-policy.hcl} :

\begin{lstlisting}[language=hcl]
path "secret/data/ci/*" {
  capabilities = ["read"]
}

path "secret/metadata/ci/*" {
  capabilities = ["read", "list"]
}
\end{lstlisting}

Jenkins peut uniquement \textbf{lire} les secrets situés sous \vpath{secret/ci/}. Il ne
peut ni les modifier, ni les supprimer, ni accéder à d'autres chemins. En cas de
compromission de Jenkins, la surface d'attaque sur Vault est donc minimale.

\subsection{Étape 4 — Le catalogue des secrets}

\begin{center}
\begin{tabularx}{\linewidth}{>{\ttfamily\small}l >{\ttfamily\small}l X}
\toprule
\textrm{\textbf{Chemin Vault}} & \textrm{\textbf{Clés}} & \textbf{Utilisé par le stage} \\
\midrule
secret/ci/dockerhub        & username, password   & Docker Push, UAT Deploy, PREPROD Deploy \\
secret/ci/nvd              & api\_key             & OWASP Dependency Check \\
secret/ci/keycloak-uat     & url, client\_secret  & UAT - Deploy \\
secret/ci/keycloak-preprod & url, client\_secret  & PREPROD - Deploy \\
\bottomrule
\end{tabularx}
\end{center}

\subsection{Étape 5 — Configuration de Jenkins}

\subsubsection*{5.a — Le plugin}

Ajout dans \vpath{ci/jenkins/plugins.txt}, installé lors du build de l'image :

\begin{lstlisting}
# --- Gestion des secrets ---
hashicorp-vault-plugin
\end{lstlisting}

\subsubsection*{5.b — Volume et variables du service Jenkins}

\begin{lstlisting}[language=yaml]
jenkins:
  volumes:
    - vault_approle:/var/jenkins_vault:ro     # role_id / secret_id AppRole
  environment:
    - VAULT_ADDR=http://vault:8200
  depends_on:
    vault-init:
      condition: service_completed_successfully
\end{lstlisting}

\subsubsection*{5.c — Script d'initialisation Groovy}

Le script \vpath{ci/jenkins/init.groovy.d/03-vault.groovy} s'exécute à chaque démarrage
de Jenkins :

\begin{lstlisting}[language=groovy]
def cred = new VaultAppRoleCredential(
    CredentialsScope.GLOBAL, 'vault-approle', 'AppRole Jenkins -> Vault',
    roleIdFile.text.trim(), Secret.fromString(secretIdFile.text.trim()), 'approle')
// ... ajout ou mise à jour dans le store global

def cfg = new VaultConfiguration()
cfg.setVaultUrl(System.getenv('VAULT_ADDR') ?: 'http://vault:8200')
cfg.setVaultCredentialId('vault-approle')
cfg.setEngineVersion(2)
GlobalConfiguration.all().get(GlobalVaultConfiguration).setConfiguration(cfg)
\end{lstlisting}

Il lit le \texttt{role\_id} et le \texttt{secret\_id} dans le volume, crée (ou met à
jour) le credential \texttt{vault-approle}, puis configure globalement le plugin (URL,
credential, KV version 2). \textbf{Aucune configuration manuelle n'est nécessaire.}

\subsection{Étape 6 — Utilisation dans le Jenkinsfile}

Des fonctions décrivent, pour chaque groupe de secrets, la correspondance entre une clé
Vault et une variable d'environnement :

\begin{lstlisting}[language=groovy]
def dockerHubSecrets() {
    return [[path: 'secret/ci/dockerhub', secretValues: [
        [envVar: 'DOCKER_USER', vaultKey: 'username'],
        [envVar: 'DOCKER_PASS', vaultKey: 'password']
    ]]]
}

def keycloakSecrets(String env) {
    return [[path: "secret/ci/keycloak-${env}", secretValues: [
        [envVar: 'KEYCLOAK_URL', vaultKey: 'url'],
        [envVar: 'KC_SECRET',    vaultKey: 'client_secret']
    ]]]
}
\end{lstlisting}

\paragraph{Avant (Jenkins Credentials)}
\begin{lstlisting}[language=groovy]
withCredentials([usernamePassword(credentialsId: 'dockerhub-credentials',
                 usernameVariable: 'DOCKER_USER', passwordVariable: 'DOCKER_PASS')]) {
    sh 'echo "${DOCKER_PASS}" | docker login -u "${DOCKER_USER}" --password-stdin'
}
\end{lstlisting}

\paragraph{Après (Vault)}
\begin{lstlisting}[language=groovy]
withVault(vaultSecrets: dockerHubSecrets() + keycloakSecrets('uat')) {
    sh 'echo "${DOCKER_PASS}" | docker login -u "${DOCKER_USER}" --password-stdin'
}
\end{lstlisting}

\begin{note}
Les fonctions sont déclarées avec \texttt{def nom() \{ \ldots \}} plutôt qu'avec des
variables \texttt{def x = \ldots} : dans un pipeline déclaratif, une variable locale
déclarée hors du bloc \texttt{pipeline \{\}} n'est pas toujours visible à l'intérieur,
alors qu'une méthode l'est toujours.
\end{note}

Les secrets n'existent dans l'environnement \textbf{que pendant l'exécution du bloc}
\texttt{withVault} et sont automatiquement masqués (\texttt{****}) dans les logs.

% =============================================================================
\section{Flux d'exécution pendant un build}
% =============================================================================

\begin{center}
\begin{tikzpicture}[
  font=\small,
  actor/.style={draw, thick, rounded corners, minimum width=2.6cm, minimum height=0.8cm,
                align=center},
  msg/.style={-{Stealth[length=2.5mm]}, thick},
  ret/.style={-{Stealth[length=2.5mm]}, thick, dashed},
  lbl/.style={font=\scriptsize, fill=white, inner sep=1pt}
]
  \node[actor, fill=green!6, draw=okgreen] (p) at (0,0)   {Pipeline};
  \node[actor, fill=blue!8,  draw=kw]      (j) at (4.5,0) {Plugin Vault\\(Jenkins)};
  \node[actor, fill=vaultyellow!70, draw=vaultdark] (v) at (9,0) {Vault};

  \foreach \n in {p,j,v} \draw[gray, dashed] (\n.south) -- ++(0,-6.4);

  \draw[msg] (0,-1.1)   -- node[lbl, above]{1. \texttt{withVault(\ldots)}} (4.5,-1.1);
  \draw[msg] (4.5,-2.0) -- node[lbl, above]{2. login AppRole (role\_id + secret\_id)} (9,-2.0);
  \draw[ret] (9,-2.9)   -- node[lbl, above]{3. token temporaire (TTL 1\,h, policy jenkins-ci)} (4.5,-2.9);
  \draw[msg] (4.5,-3.8) -- node[lbl, above]{4. lecture \texttt{secret/data/ci/dockerhub}} (9,-3.8);
  \draw[ret] (9,-4.7)   -- node[lbl, above]{5. \{username, password\}} (4.5,-4.7);
  \draw[ret] (4.5,-5.6) -- node[lbl, above]{6. variables masquées} (0,-5.6);
\end{tikzpicture}
\end{center}

\begin{enumerate}[leftmargin=*]
  \item le pipeline entre dans un bloc \texttt{withVault} ;
  \item le plugin s'authentifie auprès de Vault avec le couple \texttt{role\_id} /
        \texttt{secret\_id} ;
  \item Vault vérifie l'identité et délivre un \textbf{token temporaire} limité par la
        policy \texttt{jenkins-ci} ;
  \item le plugin lit le secret demandé avec ce token ;
  \item Vault renvoie les valeurs (uniquement si la policy l'autorise) ;
  \item les valeurs sont injectées dans les variables d'environnement, masquées dans
        les logs, puis supprimées à la sortie du bloc.
\end{enumerate}

% =============================================================================
\section{Guide d'utilisation}
% =============================================================================

\subsection{Premier démarrage}

\begin{lstlisting}[language=bash]
# 1. Préparer les secrets (fichier ignoré par Git)
cp ci/vault/secrets.env.example ci/vault/secrets.env
#    puis éditer : DOCKERHUB_TOKEN, NVD_API_KEY, KEYCLOAK_CLIENT_SECRET_*...

# 2. Construire et démarrer la stack CI
docker compose -f ci/docker-compose-ci-tools.yml up -d --build

# 3. Vérifier l'initialisation
docker logs vault-init                              # "[vault-init] Terminé"
docker logs jenkins 2>&1 | grep "\[INIT\] Vault"    # "Vault configuré → ..."
\end{lstlisting}

\subsection{Contenu de \texttt{secrets.env}}

\begin{lstlisting}
DOCKERHUB_USERNAME=moustaphisene
DOCKERHUB_TOKEN=
NVD_API_KEY=
KEYCLOAK_URL_UAT=http://host.docker.internal:8080
KEYCLOAK_CLIENT_SECRET_UAT=
KEYCLOAK_URL_PREPROD=http://host.docker.internal:8080
KEYCLOAK_CLIENT_SECRET_PREPROD=
\end{lstlisting}

\subsection{Consulter et modifier les secrets}

\paragraph{Interface web} \url{http://localhost:8200}, méthode \emph{Token}, valeur
\texttt{root}, puis \vpath{secret/} $\rightarrow$ \vpath{ci/}.

\paragraph{Ligne de commande}
\begin{lstlisting}[language=bash]
# Lire un secret
docker exec -e VAULT_TOKEN=root vault vault kv get secret/ci/dockerhub

# Modifier une seule clé (les autres sont conservées)
docker exec -e VAULT_TOKEN=root vault vault kv patch secret/ci/dockerhub password=<token>

# Historique des versions (KV v2)
docker exec -e VAULT_TOKEN=root vault vault kv metadata get secret/ci/dockerhub

# Tester les droits de Jenkins : doit échouer (lecture seule)
docker exec vault sh -c 'export VAULT_TOKEN=$(vault write -field=token \
    auth/approle/login role_id=jenkins-ci secret_id=<secret_id>); \
    vault kv put secret/ci/dockerhub password=x'
\end{lstlisting}

% =============================================================================
\section{Dépannage}
% =============================================================================

\begin{center}
\small
\begin{tabularx}{\linewidth}{>{\raggedright\arraybackslash}p{4.2cm} >{\raggedright\arraybackslash}X >{\raggedright\arraybackslash}X}
\toprule
\textbf{Symptôme} & \textbf{Cause probable} & \textbf{Solution} \\
\midrule
Jenkins ne démarre pas & \texttt{vault-init} en échec & \texttt{docker logs vault-init} \\
\texttt{Vault credentials not found for 'vault-approle'} & \texttt{03-vault.groovy} n'a pas été exécuté ou a échoué & Vérifier \texttt{vault-init}, puis \texttt{docker restart jenkins} \\
\texttt{permission denied} dans \texttt{withVault} & Chemin hors policy ou secret absent & Vérifier le chemin dans Vault et \texttt{jenkins-policy.hcl} \\
\texttt{invalid role or secret ID} & Vault redémarré (mémoire vidée) & \texttt{docker compose \ldots{} up -d} pour relancer \texttt{vault-init} \\
\texttt{init.sh: not found} ou \texttt{\$'\textbackslash r'} & Fins de ligne CRLF (Windows) & \texttt{.gitattributes} force LF ; refaire un checkout du fichier \\
Erreur sur \texttt{env\_file.required} & Docker Compose trop ancien & Mettre à jour vers Compose $\geq$ 2.24 \\
\bottomrule
\end{tabularx}
\end{center}

% =============================================================================
\section{Sécurité : mode démonstration vs production}
% =============================================================================

\begin{attention}[Mode dev — démonstration uniquement]
Le mode \texttt{-dev} stocke toutes les données \textbf{en mémoire} : un redémarrage du
conteneur \texttt{vault} efface tous les secrets (ils sont recréés par
\texttt{vault-init} depuis \texttt{secrets.env}). Le token \texttt{root} est fixe et
connu, la communication se fait en HTTP. \textbf{Ne jamais utiliser ce mode en
production.}
\end{attention}

\begin{center}
\begin{tabularx}{\linewidth}{>{\raggedright\arraybackslash}p{3.3cm} X X}
\toprule
\textbf{Aspect} & \textbf{Ce projet (mode dev)} & \textbf{Production} \\
\midrule
Stockage        & En mémoire                      & Persistant (Integrated Storage / Raft), haute disponibilité \\
Transport       & HTTP                            & TLS obligatoire (certificats) \\
Descellement    & Automatique (dev)               & Clés de Shamir ou auto-unseal (AWS KMS, Azure Key Vault\ldots) \\
Accès admin     & Token \texttt{root} fixe        & Token root révoqué ; humains via OIDC/LDAP + MFA \\
\texttt{secret\_id} & Permanent (\texttt{secret\_id\_ttl=0}) & Courte durée, usage limité, \emph{response wrapping} \\
Audit           & Non activé                      & \texttt{vault audit enable file} $\rightarrow$ SIEM \\
Secrets         & Statiques                       & Dynamiques quand c'est possible (bases de données, cloud) \\
\bottomrule
\end{tabularx}
\end{center}

\subsection{Bénéfices obtenus}

\begin{itemize}[leftmargin=*]
  \item[\textcolor{okgreen}{\checkmark}] Plus aucun secret saisi dans Jenkins ni versionné dans Git ;
  \item[\textcolor{okgreen}{\checkmark}] Authentification machine (AppRole) avec tokens à durée de vie limitée ;
  \item[\textcolor{okgreen}{\checkmark}] Moindre privilège : Jenkins en lecture seule sur \vpath{secret/ci/*} ;
  \item[\textcolor{okgreen}{\checkmark}] Secrets masqués dans les logs et limités à la portée du bloc \texttt{withVault} ;
  \item[\textcolor{okgreen}{\checkmark}] Versionnage des secrets (KV v2) et rotation centralisée ;
  \item[\textcolor{okgreen}{\checkmark}] Configuration 100\,\% automatisée et reproductible.
\end{itemize}

\subsection{Pistes d'amélioration}

\begin{enumerate}[leftmargin=*]
  \item Passer Vault en mode serveur avec stockage Raft, TLS et auto-unseal ;
  \item Activer l'\textbf{audit device} pour tracer chaque lecture de secret ;
  \item Utiliser des \textbf{secrets dynamiques} (moteur \texttt{database}) pour les
        bases PostgreSQL ;
  \item Faire lire le secret client Keycloak directement par les applications Spring
        Boot via \textbf{Spring Cloud Vault} ;
  \item Gérer la configuration de Vault en \emph{Infrastructure as Code} avec Terraform.
\end{enumerate}

% =============================================================================
\appendix
\section{Aide-mémoire des commandes}
% =============================================================================

\begin{center}
\small
\begin{tabularx}{\linewidth}{>{\ttfamily\raggedright\arraybackslash}p{8.6cm} X}
\toprule
\textrm{\textbf{Commande}} & \textbf{Effet} \\
\midrule
vault status                               & État du serveur (scellé, version\ldots) \\
vault auth list                            & Méthodes d'authentification activées \\
vault policy read jenkins-ci               & Afficher la policy Jenkins \\
vault read auth/approle/role/jenkins-ci    & Paramètres du rôle AppRole \\
vault kv list secret/ci                    & Lister les secrets du pipeline \\
vault kv get secret/ci/nvd                 & Lire un secret \\
vault kv put secret/ci/nvd api\_key=\ldots  & Créer ou remplacer un secret \\
vault kv patch secret/ci/nvd api\_key=\ldots & Modifier une clé \\
vault kv rollback -version=1 secret/ci/nvd & Restaurer une version précédente \\
vault token lookup                         & Informations sur le token courant \\
\bottomrule
\end{tabularx}
\end{center}

\noindent Dans le conteneur : préfixer par
\texttt{docker exec -e VAULT\_TOKEN=root vault}.

\end{document}
